% เรียงความ 1 หน้า A4: เขียนเนื้อหาตรงนี้ (ย่อหน้าใหม่ = เว้นบรรทัดว่าง)
\textbf{เครือข่ายที่ผมรู้จักในชีวิตประจำวัน}

เครือข่ายแรกที่ผมใช้ทุกวันคือเครือข่าย Wi-Fi และอินเทอร์เน็ตในหอพัก อินเทอร์เน็ตของหอมาจากสาย LAN ที่ผนังห้อง ซึ่งต่อไปยังเราเตอร์ MikroTik ของหอ ผมต่อสายนี้เข้ากับเราเตอร์ Wi-Fi 6 ของตัวเอง เพื่อปล่อยสัญญาณให้โทรศัพท์และโน้ตบุ๊กในห้อง เวลาผมเปิดเว็บไซต์ ข้อมูลจะวิ่งจากอุปกรณ์ผ่านคลื่น Wi-Fi ไปที่เราเตอร์ในห้อง ต่อไปยังเครือข่ายของหอ แล้วจึงออกสู่อินเทอร์เน็ต ระหว่างทางมีมาตรฐานหลายตัวทำงานร่วมกัน เช่น IEEE 802.11ax สำหรับ Wi-Fi, DHCP ที่แจก IP Address ให้อุปกรณ์อัตโนมัติ, DNS ที่แปลงชื่อเว็บไซต์เป็น IP และ HTTPS ที่เข้ารหัสข้อมูลระหว่างผมกับเว็บไซต์

ปัญหาที่ผมเจอคือหอพักใช้ระบบ Hotspot ที่ต้องล็อกอินผ่านหน้าเว็บก่อนใช้งาน และเราเตอร์เดิมในห้องทำให้สัญญาณกระตุกและหลุดบ่อย ผมแก้ด้วยการเปลี่ยนมาใช้เราเตอร์ของตัวเอง และตั้งค่าให้ล็อกอิน Hotspot อัตโนมัติ หลังจากนั้นการเชื่อมต่อก็เสถียรขึ้นมาก

อีกเครือข่ายหนึ่งที่ผมใช้คือ Tailscale เนื่องจากเครือข่ายหอพักอยู่หลัง NAT ผมจึงตั้ง Port Forwarding ไม่ได้ ทำให้เข้าถึงเครื่องคอมพิวเตอร์ของตัวเองจากข้างนอกไม่ได้ Tailscale แก้ปัญหานี้ด้วยการสร้างเครือข่ายส่วนตัวเสมือน (Mesh VPN) ทุกเครื่องที่อยู่ในเครือข่ายจะได้ IP ส่วนตัวในช่วง 100.x.x.x และสื่อสารกันได้เหมือนอยู่ใน LAN เดียวกัน ไม่ว่าจะอยู่ที่ไหน หลักการคือมีเซิร์ฟเวอร์กลางทำหน้าที่แจกกุญแจและบอกว่าแต่ละเครื่องอยู่ที่ไหน ส่วนข้อมูลจริงจะส่งตรงระหว่างเครื่องแบบ Peer-to-Peer โดยใช้เทคนิค NAT Traversal และเข้ารหัสด้วยโปรโตคอล WireGuard หากเชื่อมต่อตรงไม่ได้ ข้อมูลจะถูกส่งอ้อมผ่านเซิร์ฟเวอร์ Relay ที่เรียกว่า DERP แทน เครือข่ายนี้ทำให้ผมเข้าถึงเครื่องที่หอได้จากทุกที่อย่างปลอดภัย

\textbf{เครือข่ายที่ผมฝันถึง}

เครือข่ายที่ผมฝันถึงคือเครือข่ายที่ไม่มีกำแพงกั้นระหว่างอุปกรณ์ จากประสบการณ์ที่ต้องใช้ Tailscale เพื่อหลบข้อจำกัดของ NAT ผมอยากเห็นเครือข่ายที่ใช้ IPv6 อย่างเต็มรูปแบบ ทุกอุปกรณ์จะมี IP Address สาธารณะเป็นของตัวเอง และเชื่อมต่อถึงกันได้โดยตรงตามหลัก End-to-End ที่ออกแบบอินเทอร์เน็ตไว้ตั้งแต่แรก นักศึกษาจะเปิดเซิร์ฟเวอร์ทดลองหรือเข้าถึงเครื่องของตัวเองจากที่ไหนก็ได้ โดยไม่ต้องพึ่งเครื่องมือเสริม แต่การเปิดให้เชื่อมต่อตรงได้ก็ต้องมี Firewall ที่ผู้ใช้ตั้งค่าเองได้ง่ายควบคู่กันไปด้วย เพื่อเลือกได้ว่าจะเปิดให้ใครเข้าถึงบ้าง

นอกจากนี้ ผมอยากให้เครือข่ายทั้งมหาวิทยาลัยเป็นเครือข่ายเดียวกัน ไม่ต้องล็อกอินผ่านหน้าเว็บทุกครั้ง แต่ยืนยันตัวตนด้วยบัญชีนักศึกษาเพียงครั้งเดียว คล้ายกับระบบ eduroam เมื่อเดินจากหอพักไปตึกเรียน อุปกรณ์จะสลับไปใช้จุดกระจายสัญญาณที่ใกล้ที่สุดเองโดยที่การเชื่อมต่อไม่หลุด ซึ่งทำได้ด้วยมาตรฐานอย่าง IEEE 802.11r ที่ช่วยให้ย้ายจุดเชื่อมต่อได้เร็ว

สุดท้าย ผมอยากให้เครือข่ายดูแลตัวเองได้ เมื่อมีจุดไหนสัญญาณรบกวนหรือมีผู้ใช้หนาแน่น ระบบจะตรวจพบและปรับช่องสัญญาณหรือกระจายผู้ใช้ไปยังจุดอื่นอัตโนมัติ แทนที่ผู้ใช้จะต้องแก้ปัญหาเองเหมือนที่ผมเคยทำ ผู้ใช้จะไม่ต้องสนใจว่าเบื้องหลังเครือข่ายทำงานอย่างไร เพราะมัน "แค่ใช้งานได้"
